\documentclass[10pt,letterpaper]{article}
\usepackage{cornell-notes}

\title{Clause 02: Normative References}
\author{}
\date{}

\setDocTitle{Clause 02: Normative References}
\setDocSubtitle{{C 2024}}
\setDocOwner{{C 2024 Cornell Notes}}
\setDocDate{{\today}}
\setCornellCollection{{C 2024 Cornell Notes}}
\setCornellUnitType{{Clause}}
\setCornellUnitNumber{02}
\setCornellUnitTitle{Normative References}
\hypersetup{pdftitle={Clause 02: Normative References},pdfsubject={C 2024 study notes},pdfauthor={}}

\begin{document}
\maketitle



\begin{CornellOverviewBox}{Clause Orientation}
This clause identifies external specifications whose cited content becomes part of the requirements when invoked by the source material. It also defines how dated and undated references select an applicable edition. The references supply terminology, character data, mathematical conventions, date/time and currency representations, and floating-point rules that the remaining clauses depend on.
\end{CornellOverviewBox}

\section{Learning Objectives}
\begin{itemize}
  \item Explain when an external reference contributes normative requirements.
  \item Distinguish dated references from undated references.
  \item Identify the technical dependency categories supplied by referenced documents.
  \item Avoid treating every bibliography entry or external resource as normative.
  \item Trace a requirement through a reference without losing the calling clause context.
\end{itemize}

\section{Normative Structure Map}
\begin{CornellNotesTable}
\CornellNoteRow{Dated reference}{Only the specifically cited edition applies for that dependency.}
\CornellNoteRow{Undated reference}{The latest applicable edition, including amendments, supplies the cited material.}
\CornellNoteRow{Normative force}{A cited document contributes requirements only where the calling clause invokes it.}
\end{CornellNotesTable}

\section{How References Become Normative}
\begin{CornellNotesTable}
\CornellNoteRow{What makes a reference normative?}{Its content is cited in a way that makes some or all of that content necessary to satisfy requirements in the source material.}
\CornellNoteRow{Does listing a document import everything it says?}{Not automatically. The invoking text determines which referenced content matters and how it participates in the requirement.}
\CornellNoteRow{How is this different from a bibliography?}{A bibliography supports background or further reading. A normative reference can supply definitions or rules required to interpret or implement a provision.}
\CornellNoteRow{What is the safe reading method?}{Read the calling clause, identify the exact dependency, then consult the selected edition of the referenced material. Do not detach imported wording from its use context.}
\end{CornellNotesTable}

\begin{CornellSummaryBox}{Section Summary}
Normativity arises from the way a reference is used, not merely from the existence of a citation.
\end{CornellSummaryBox}



\section{Dated and Undated References}
\begin{CornellNotesTable}
\CornellNoteRow{What does a dated reference select?}{Only the specifically cited edition applies for that dependency.}
\CornellNoteRow{What does an undated reference select?}{The latest applicable edition, including amendments, supplies the referenced content.}
\CornellNoteRow{Why does the distinction matter?}{Definitions, character properties, arithmetic details, or interchange formats can evolve. Selecting the wrong edition can change the interpreted requirement.}
\CornellNoteRow{What should an implementation record?}{The edition or data version actually used for each externally supplied dependency, especially where the reference is undated and can change over time.}
\end{CornellNotesTable}

\begin{CornellSummaryBox}{Section Summary}
Reference versioning is part of conformance evidence: dated citations are fixed, while undated citations follow the latest applicable material.
\end{CornellSummaryBox}



\section{Dependency Categories}
\begin{CornellNotesTable}
\CornellNoteRow{Which terminology dependency is supplied?}{A foundational information-technology vocabulary and mathematical notation conventions support consistent definitions.}
\CornellNoteRow{Which representation dependencies are supplied?}{External sources define universal character repertoires and properties, currency codes, and date/time interchange conventions used by related language or library facilities.}
\CornellNoteRow{Which arithmetic dependency is supplied?}{A referenced floating-point specification supplies arithmetic models and behavior for implementations that support the associated features.}
\CornellNoteRow{How should online character-property data be treated?}{Use the referenced dataset and its selected version exactly where the language rules depend on derived character properties; do not substitute an unrelated database.}
\end{CornellNotesTable}

\begin{CornellSummaryBox}{Section Summary}
The references are functional dependencies: terminology, encoding, arithmetic, and interchange provisions are imported only where later clauses call on them.
\end{CornellSummaryBox}



\section{Programmer and Implementation Checklist}
\begin{itemize}
  \item[$\square$] Locate the exact normative sentence that invokes an external reference.
  \item[$\square$] Determine whether the reference is dated or undated before selecting an edition.
  \item[$\square$] Record the selected edition or data version as part of implementation evidence.
  \item[$\square$] Apply only the referenced content needed by the calling provision.
  \item[$\square$] Do not upgrade a dated reference silently.
  \item[$\square$] Do not treat informative background material as a conformance requirement.
  \item[$\square$] Recheck undated online datasets when validating a current implementation.
\end{itemize}

\begin{CornellSummaryBox}{Clause Synthesis}
Clause 2 makes the specification a controlled composition of internal rules and externally supplied definitions. Correct use requires both dependency tracing and version discipline. A referenced document is not a free-standing replacement for the calling clause: the language rule determines what is imported, while the dated-or-undated form determines which version supplies it.
\end{CornellSummaryBox}



\clearpage
\section{Self-Test Questions}
\begin{enumerate}
  \item When does an external document contribute requirements?
  \item Which edition applies to a dated reference?
  \item Which edition applies to an undated reference?
  \item Why is a bibliography entry not automatically normative?
  \item An implementation uses a newer edition for a dependency that was cited by date. Is that automatically conforming?
  \item A character-property rule depends on an undated online dataset. What should a validation record contain?
\end{enumerate}

\clearpage
\section{Answer Key}
\begin{enumerate}
  \item When the source refers to its content in a way that makes that content necessary to interpret or satisfy a requirement.

  \item Only the edition explicitly cited.

  \item The latest applicable edition, including applicable amendments.

  \item Normative force depends on incorporation by a requirement. A bibliography can be informative and need not constrain conformance.

  \item No. The cited edition remains controlling unless the source permits or adopts the newer one.

  \item It should identify the dataset version used and show that the relevant property was interpreted in the context of the calling language rule.
\end{enumerate}

\section{Key-Term Recap}
\begin{CornellNotesTable}
\toprule
Term & Working meaning \\
\midrule
\endfirsthead
\toprule
Term & Working meaning (continued) \\
\midrule
\endhead
Normative reference & An external source whose cited content participates in one or more requirements. \\
Dated reference & A reference for which only the specifically cited edition applies. \\
Undated reference & A reference resolved to the latest applicable edition, including amendments. \\
Calling clause & The internal provision that invokes and gives context to referenced content. \\
Dependency trace & Evidence connecting an internal requirement to the relevant external definition or rule. \\
\bottomrule
\end{CornellNotesTable}

\begin{CornellSummaryBox}{One-Sentence Takeaway}
A referenced source affects conformance only through the provision that invokes it, using the edition-selection rule implied by whether the citation is dated or undated.
\end{CornellSummaryBox}
\end{document}
